% =====================================================================
% ESQUELETO — ERCEMAPI 2026
% Baixe o template oficial em:
% https://www.sbc.org.br/wp-content/uploads/2024/07/modelosparapublicaodeartigos.zip
% e substitua o preambulo abaixo pelo sbc-template.cls oficial deles.
% Este arquivo assume os comandos padrao do sbc-template (\title, \author,
% \address, ambientes resumo/abstract) — ajuste se a versao baixada diferir.
% =====================================================================
\documentclass[12pt]{article}
\usepackage{sbc-template}
\usepackage[brazil]{babel}
\usepackage[utf8]{inputenc}
\usepackage{graphicx}
\usepackage{booktabs}

% ---------------------------------------------------------------------
% DOUBLE BLIND: nao preencher \author nem \address nesta fase.
% O proprio template SBC tem um modo de ocultar essa secao — checar no
% pacote baixado. Nao deixar nomes de autores, IFCE, LISA, nem "Luiz
% Henrique" em lugar nenhum do corpo do texto ou dos agradecimentos.
% ---------------------------------------------------------------------
\title{Avaliação Exploratória de um Modelo Compacto de \emph{Tool-Calling}
para Interfaces de Comando em Português no Contexto de IoT Apícola}

% \author{Omitido para revisão double-blind}
% \address{Omitido para revisão double-blind}

\begin{document}

\maketitle

\begin{resumo}
% TODO (~150 palavras): objetivo, metodo (motor host + hardware real
% ESP32-S3 N16R8), 36 casos de teste, principais metricas encontradas,
% uma frase de conclusao/contribuicao.
\end{resumo}

\begin{abstract}
% TODO: versao em ingles do resumo acima.
\end{abstract}

% =====================================================================
\section{Introdução}
% TODO, ~1 pagina. Sugestao de estrutura:
% - Contexto: IoT aplicado a apicultura em regiao de semiarido, operadores
%   nem sempre tecnicos, conectividade instavel em campo.
% - Problema: interfaces de comando tradicionais (comandos seriais fixos,
%   paineis web) exigem que o operador memorize sintaxe.
% - Lacuna: modelos de linguagem convencionais sao pesados demais para
%   microcontroladores; a maioria dos benchmarks de tool-calling compacto
%   e feita em ingles, sem avaliacao em portugues nem em dominio agricola.
% - Contribuicao deste trabalho (deixar claro que e EXPLORATORIO,
%   sem validacao em campo com colmeia real):
%   1) schema de ferramentas para comando por linguagem natural no
%      dominio apicola;
%   2) conjunto de teste de 36 frases em portugues cobrindo comandos
%      diretos, variacoes coloquiais, parametros implicitos, ambiguidade
%      e frases fora de escopo;
%   3) avaliacao end-to-end no hardware-alvo real (ESP32-S3 N16R8),
%      incluindo latencia por tipo de comando.

% =====================================================================
\section{Trabalhos Relacionados}
% TODO, ~1 pagina. Citar:
% - Needle 2 / Cactus Compute (arXiv:2607.18363) — arquitetura Simple
%   Attention Network, quantizacao Cactus Quants CQ2-bit, resultados em
%   Mobile Actions / DroidCall / Seal-Tools / BFCL v4.
% - Trabalhos de IoT em apicultura (Colmeia v2.0 e correlatos da
%   literatura — buscar 2-4 referencias de monitoramento de colmeias
%   com sensores, sem citar o proprio projeto de forma identificavel).
% - Trabalhos de interfaces de comando por voz/texto em dispositivos de
%   borda (ESP-SR da Espressif, outros modelos pequenos de tool-calling
%   como FunctionGemma, LFM2.5, mencionados no material do Cactus).
% - Posicionar a lacuna: nenhum desses avalia comando em portugues nem
%   dominio agricola regional.

% =====================================================================
\section{Proposta}
% TODO, ~1-1.5 pagina. Descrever a arquitetura textualmente e referenciar
% a Figura 1 (fluxograma). Detalhar:
% - As 3 ferramentas do schema (consultar_status_colmeia,
%   listar_colmeias_alerta, ativar_modo_inspecao) e a decisao de manter
%   o schema pequeno (orcamento de ~256 tokens de contexto do Needle 2).
% - O fluxo: comando em texto -> inferencia local restrita por gramatica
%   -> chamada valida ou chamada vazia -> execucao do handler -> resposta.
% - Deixar explicito que os handlers usam dados MOCK nesta fase (nao ha
%   integracao real com sensores de uma estacao Colmeia neste experimento).

\begin{figure}[h]
  \centering
  % Exportar o SVG do fluxograma como PDF antes de incluir (Inkscape,
  % ou "Export as PDF" em qualquer visualizador de SVG).
  \includegraphics[width=0.85\textwidth]{fluxo_needle2_colmeia.pdf}
  \caption{Fluxo de funcionamento do comando por linguagem natural.}
  \label{fig:fluxo}
\end{figure}

% =====================================================================
\section{Metodologia}
% TODO, ~1.5 pagina. Subsecoes sugeridas:

\subsection{Hardware e ambiente}
% TODO: ESP32-S3-WROOM-1 N16R8 (16 MB flash, 8 MB PSRAM octal), motor
% host para ensaios preliminares, ESP-IDF v5.5+, modelo needle2.cact
% (Cactus Compute, Apache 2.0).

\subsection{Schema de ferramentas}
% TODO: reproduzir o schema compactado (ou referenciar apendice), com
% justificativa do numero de ferramentas dado o limite de contexto.

\subsection{Conjunto de teste}
% TODO: descrever as 36 frases e as 5 categorias (diretos, variacoes de
% fala, parametro implicito, ambiguos, fora de escopo). Justificar que
% o conjunto foi construido manualmente para cobrir casos de uso reais
% de um apicultor, incluindo casos-limite deliberados.

\subsection{Protocolo experimental e metricas}
% TODO: descrever a execucao automatizada (script sobre a TUI de
% referencia), numero de repeticoes por frase (se houve), e definir
% formalmente as 4 metricas:
% - acuracia de correspondencia exata
% - taxa de bem-formado (JSON valido contra o schema)
% - taxa de rejeicao correta (fora_escopo -> chamada vazia)
% - taxa de omissao correta de parametro nao fundamentado
% - latencia media/minima/maxima por requisicao

% =====================================================================
\section{Resultados}
% TODO apos rodar o experimento. Tabelas sugeridas:

\begin{table}[h]
\centering
\caption{Métricas agregadas por categoria de comando.}
\label{tab:metricas}
\begin{tabular}{lccc}
\toprule
Categoria & N & Acurácia exata & Latência média (s) \\
\midrule
Diretos            & 8 & TODO & TODO \\
Variações de fala   & 8 & TODO & TODO \\
Parâmetro implícito & 6 & TODO & TODO \\
Ambíguos            & 6 & TODO & TODO \\
Fora de escopo      & 8 & TODO & TODO \\
\bottomrule
\end{tabular}
\end{table}

% TODO: uma segunda tabela ou paragrafo discutindo os casos individuais
% mais interessantes (principalmente das categorias ambiguos e
% parametro_implicito, que sao ricos em nuance).

% =====================================================================
\section{Discussão}
% TODO, ~1 pagina. Pontos a cobrir:
% - O modelo foi treinado em ingles; discutir o impacto disso nos
%   resultados em portugues, com base nos numeros da Secao 5.
% - Viabilidade real de embarque no ESP32-S3 N16R8 (memoria, latencia
%   medida vs. os ~25-47s reportados pelo Cactus em outros dispositivos).
% - Quando faz sentido usar um modelo vs. regras deterministicas (conectar
%   com a decisao de projeto do firmware da estacao meteorologica, que
%   usa regras fixas por restricao de memoria do ESP32 classico).
% - Limitacoes: dados mock, sem validacao com apicultor real, dataset
%   pequeno (36 casos), sem fine-tuning.

% =====================================================================
\section{Trabalhos Futuros}
% TODO. Incluir pelo menos:
% - Integracao com a API real da estacao de monitoramento de colmeias.
% - "O hardware utilizado ja dispoe de microfone integrado; a integracao
%   de um pipeline de reconhecimento de fala em portugues, compativel
%   com o orcamento de memoria do dispositivo, e apontada como extensao
%   natural deste trabalho."
% - Fine-tuning do Needle 2 com corpus proprio de comandos apicolas em
%   portugues (processo local, documentado pelo proprio projeto Cactus).
% - Validacao de campo com apicultores reais.

% =====================================================================
\section{Conclusão}
% TODO, ~0.5 pagina. Retomar objetivo, principal achado quantitativo,
% e a mensagem central (viabilidade + limitacoes claras + caminho futuro).

% =====================================================================
\bibliographystyle{sbc}
\bibliography{referencias}
% TODO: criar referencias.bib com pelo menos:
% - Cactus Compute / Needle 2 (arXiv:2607.18363)
% - 2-4 referencias de IoT em apicultura
% - 1-2 referencias de tool-calling / modelos compactos para edge
% - Documentacao ESP-IDF / ESP32-S3, se citada tecnicamente

\end{document}
